% One-time imported draft. Editable native source; no Markdown regeneration.
\section{An address identifies a
  location}\label{001-console-calculation-p04-an-address-identifies-a-location}

Loading the prompt address puts a location into \texttt{\$a0}. The
string remains in memory. This arrangement lets a service read a message
of several bytes using one register-sized argument: the argument
identifies where the message begins.

A useful notation is \texttt{A} for the address of the first byte. The
next byte has address \texttt{A\ +\ 1}, and the one after that has
address \texttt{A\ +\ 2}. Memory is byte-addressed. The numerical
address is a property of the assembled layout, so this lesson refers to
the label instead of relying on one fixed hexadecimal address.

\texttt{la} is an assembler convenience called a pseudo-instruction. Its
expansion may use more than one machine instruction. A source-level
operation and an executed machine instruction therefore need not be the
same unit when stepping or counting work.
